\section{Discussion}
\label{sec:limitations}

The experiments separate capabilities that are often presented together. The Craftax checkpoint
supports nearby visual-goal planning, while its reward-planning pipeline has weak held-out reward and value
predictions. The gate encoders retain useful obstacle information under an added decoder and CNN,
while their tested imagination pipelines fail to discover the strategy. Nuisance cubes and the
population constructions show why a marginal non-collapse condition alone cannot settle either
question.

The practical consequence is to evaluate a world model at each interface: information recoverable
from real encodings, prediction of controlled counterfactuals, head generalization, and behavior in
the real environment. A controller's failure should lead to this decomposition before it is used to
judge the representation; a successful probe should be followed by prediction and control tests
before it is used to certify imagination.

Our empirical scope is one jointly trained squared-loss objective, small visual domains, and short
training budgets. Craftax and the flow reference each use one fitted model. Gate predictors, heads,
and PPO protocols have not been isolated by a fully matched intervention. The theory gives
population existence results, with measurable constructions rather than a finite-network or
optimization analysis. In particular, the fresh-bit example does not explain loss of a predictable
wall. What the evidence does establish is that non-collapse, decodability, and useful imagined
control are separate properties, requiring separate support.
